
\section{\lr{Bamboo}}

\lr{Bamboo} یک کارگزار مجتمع سازی پیوسته\LTRfootnote{Continuous Integration} یا
\lr{CI} است که می‌تواند برای توسعه سیستم‌های نرم‌افزاری و ساخت آنها بر اساس
قراردادهای این روش مورد استفاده قرار گیرد.

% What does this mean?
مجتمع سازی پیوست یک متدلوزی توسعه و ساخت سیستم‌های نرم‌افزاری است که در آن 
هم زمان با تغییر کد منبع سیستم‌های نرم‌افزاری، فرآیند ترجمه، ارزیابی
برنامه‌ها و ارزیابی مجتمع سازی اجرا (و یا تحریک) می‌شود تا از درست بودن کدهای
ایجاد شده در مقایسه با کل سیستم اطمینان حاصل شود. در این فرآیند توسعه همواره
کیفیت نرم‌افزار ایجاد شده در حالت کلی و ایجاد بسته‌های نرم‌افزاری در قبال هر
تغییر بررسی خواهد شد. شکل \ref{image/standard/build/bamboo-arch} فرآیند کلی مورد
استفاده در این متدلوژی را نمایش داده است.

فرآیند ایجاد سنتی سیستم‌های نرم‌افزاری شامل ترجمه، ارزیابی برنامه‌ها و ایجاد
بسته‌ها همواره در این متدلوژی اجرا می‌شود اما به این معنی نیست که به واسطه هر
تغییر یک نسخه جدید از نرم‌افزار ایجاد شده است. از این رو فرآیند ایجاد نسخه جدید
از یک سیستم نرم‌افزاری کمی با حالت سنتی آن متفاوت است. در این متدلوژی ارائه یک
نسخه جدید تنها زمانی انجام می‌شود که تمام نقشه راه در نظر گرفته شده برای سیستم
نرم‌افزاری به صورت کامل پیاده سازی شده و تمام ارزیابی‌ها را با موفقیت پشت
سر گذاشته باشد.

\begin{figure}
\centering
\includegraphics[width=0.75\textwidth]{image/standard/build/bamboo-arch}
\caption[فرآیند مجتمع سازی پیوسته یا \lr{CI}]{
	در این فرآيند
}
\label{image/standard/build/bamboo-arch}
\end{figure}


% How does Bamboo do this?

% Bamboo is the central management server which schedules and coordinates all work.
اما بامبو یک کارگزار مرکزی است که تمام فرآیندهای ساخت و ارزیابی بسته‌های
نرم‌افزاری ایجاد شده را مدیریت و برنامه ریزی می‌کند. این کارگزار با استفاده از
زیر سیستم‌های موجود، نرم‌افزارهای متفاوت را ترجمه کرده و ارزیابی‌های تعریف شده
برای آنها را اجرا می‌کند.

% Bamboo itself has interfaces and plugins for lots of types of work.
گرچه افزونه‌های متفاوتی برای این کارگزار تعریف شده است که قادر هستند کارهای
متفاوتی را در فرآیند ساخت و ارزیابی اجرا کنند اما امکان تعریف افزونه‌های جدید
نیز پیش بینی شده است. توانایی تعریف افزونه‌های جدید برای این کارگزار، آن را به
عنوان یک ابزار بسیار قدرتمند در متدلوژي مجتمع سازی پیوسته تبدیل کرده است.

%  Bamboo first gets your source from a source repository (lots of plugins here for a variety of systems).
فرآیند کلی ایجاد و ارزیابی سیستم‌های نرم‌افزاری با دریافت برنامه‌های ایجاد شده
از سرورهای کنترل نسخه آغاز می‌شود. این کارگزار با استفاده از افزونه‌های متفاوت
توانایی دریافت برنامه‌های منبع از کارگزارهای متفاوت کنترل نسخه مانند \lr{Git}،
\lr{SVN}، \lr{CVS} و غیر را دارد.

% Then Bamboo starts the build - that can be done by calling something like MSBuild to build your Visual Studio solution, or Maven to call your compiler and linker - whatever you use.
فرآیند ساخت در این کارگزار با فراخوانی ابزارها معمول ساخت برنامه‌های نرم‌افزاری
اینجام می‌شود. برای نمونه با استفاده از فراخوانی مترجم \lr{GCC} در لینوکس
نرم‌افزارهایی که به زبان برنامه سازی \lr{C/C++} ایجاد شده‌اند ترجمه می‌شوند. از
این رو زیر سیستم‌های تعریف شده برای این کارگزار می‌بایست توانایی استفاده از این
فراخوانی‌های در آنها ایجاد شده باشد (به این معنی  نه تنها ابزارهای مورد نیاز
باید نصب شده باشد بلکه در مسیرهای پیش فرض سیستم قرار گرفته شده باشد).

% Once your solution or project is built, you have "artifacts" (build results, for example, an executable app, config files, etc.).
پس از اتمام فرآیند ساخت گردایه‌ای از برنامه‌های و پرونده‌هایی ایجاد خواهد شد که
به عنوان نتیجه ساخت در نظر گرفته می‌شود. نتیجه‌های ایجاد شده در فرآیند ساخت را
اصطلاحا مصنوع\LTRfootnote{Artifact} گفته می‌شود. مصنوعات ایجاد شده در فرآیند
ساخت ممکن است از یک برنامه‌اجرایی تا یک پرونده پیکره بندی متفاوت باشد. برای
نمونه مستند فنی که نتیجه فرآیند ساخت مستند بر اساس کد فنی است شامل پرونده‌های
زنگام\LTRfootnote{HTML}، تصاویر و پوشه‌های متفاوت است.

% You can do additional things with the build artifacts:
%         zip them up into a ZIP file and copy them somewhere.
%         run a install builder on them and create an MSI.
%         install them on a test server to make sure everything installs just fine.
کارگزار ساخت می‌تواند کارهای متفاوتی را بر اساس مصنوعات انجام دهد. فشرده سازی،
نصب، اجرا برنامه‌های ارزیاب بخشی از فرآیندهایی است که می‌توان بر اساس مصنوعات
اجرا کرد. برای نمونه با نصب یک برنامه ایجاد شده به روی ماشین‌های مجازی در نظر
گرفته شده، درست بودن فرآیند نصب مورد بررسی قرار می‌گیرد.

% Bamboo provides a web front-end for configuration and for reporting the status of builds.
استفاده از یک تارافزار به عنوان واسط کاربر بامبو را به یک کارگزار بسیار ساده
تبدیل کرده است که با وجود سادگی ظاهر توانایی اجرای فرآیندهای ساخت بسیار پیچیده
را دارد. با استفاده از این واسط کاربر به سادگی می‌توان نتایج به دست آمده از
فرآیندهای ساخت متفاوت را بررسی و تنظیم‌های مورد نیاز را در پروژه‌های ایجاد کرد.
 
% What does Bamboo need?
کارگزار بامبو با زمانبندی فرآیند ساخت هر نرم‌افزار، فرآیندهای ارزیابی را به روی
نتایج به دست آمده از آن را اجرا خواهد کرد. اما برای استفاده از تمام امکانهای
تعریف شده در این کارگزار نیاز است که پیش نیازهای آن فراهم شود. تمام نیازهای این
کارگزار در این فرآیند به عبارت است از:

\begin{itemize}
  \item کارگزار کنترل نسخه و کد منبع
  \item نپشته‌های ایجاد
  \item برنامه‌های ارزیابی
  \item نپشته‌های بسته بندی
  \item نپشته‌های ارزیابی مجتمع سازی
\end{itemize}

در متدلوژی مجتمع سازی پیوست فرض بر این است فردی که کد منبع را تغییر می‌دهد
مسئولیت تمام تغییرهایی خواهد بود که در کل سیستم رخ می‌دهد. این فرد می‌بایست تمام
اشکالهای احتمالی که ممکن است در کل سیستم رخ دهد را برطرف کرده تا فرآیند ساخت به
صورت کامل اجرا شود و مصنوعات ایجاد شده ارزیابی‌های تعریف شده را با موفقیت طی
کند.

نکته‌ای که باید به آن توجه کرد این است که زیر سیستم‌های مورد استفاده در فرآیند
ساخت و ارزیابی توانایی‌های خاص خود را دارند از این رو برای ایجاد نپشته‌های ساخت
و ارزیابی می‌بایست توانایی‌های مورد نیاز به درستی تعیین شود. از این رو کارگزار
تنها زمانی که زیر سیستم مناسب برای اجرای این فرایند را در اختیار داشته باشد به
اجرای نپشته‌های ساخت و ارزیابی خواهد پرداخت.

% How is a Bamboo workflow organised?

کارگزار بامبو، وظایف مورد نیاز در فرآیند مجتمع سازی را با استفاده از موجودیت‌های
متفاوتی مانند پروژه، نقشه، کار و وظیفه سازماندهی می‌کند. ساختار سلسله مراتبی این
موجودیت ها در شکل \ref{image/standard/build/bamboo-workflow-objects} نمایش داده
شده است.

در فرآیند ایجاد یک پروژه ابتدا کارگزار نخستن نقشه را انتخاب می‌کند و به اجرا آن
می‌پردازد. پس از پایان یافتن هر نقشه، نقشه بعد برای اجرا انتخاب خواهد شده در
حالی که هر نقشه نیز خود از چندین وهله تشکیل شده است. از این رو برای اجرا هر وهله
نیاز است تمام کارهای تعریف شده در آن به ترتیب اجرا شود. برای اجرای هر کار نیز
نیاز است که وظایف تعریف شده در آن به ترتیب اجرا شود. بنابر این برای اجرای هر
نقشه از یک پروژه نیاز است تمام وظیفه‌ها، کارها و وهله‌ها به صورت کامل اجرا شوند
تا آن نقشه تکمیل شود.


\begin{figure}
\centering
\includegraphics[width=0.75\textwidth]{image/standard/build/bamboo-workflow-objects}
\caption[فرآیند مجتمع سازی پیوسته یا \lr{CI}]{
	در این فرآيند
}
\label{image/standard/build/bamboo-workflow-objects}
\end{figure}


% Project	
% 
%     Has one, or more, plans.
%     Provides reporting (using the wallboard, for example) across all plans in the project.
%     Provides links to other applications.
% 
% Plan	
% 
%     Has a single stage, by default, but can be used to group jobs into multiple stages.
%     Processes a series of one or more stages that are run sequentially using the same repository.
%     Specifies the default repository.
%     Specifies how the build is triggered and triggering dependencies between a plan and other plans in the project.
%     Specifies notifications of build results.
%     Specifies who has permission to view and configure the plan and its jobs.
%     Provides for the definition of plan variables.
% 
% Stage	
% 
%     Has a single job, by default, but can be used to group multiple jobs.
%     Processes its jobs in parallel, on multiple agents (where available).
%     Must successfully complete all its jobs before the next stage in the plan can be processed.
%     May produce artifacts that can be made available for use by a subsequent stage.
% 
% Job	
% 
%     Processes a series of one or more tasks that are run sequentially on the same agent.
%     Controls the order in which tasks are performed.
%     Collects the requirements of individual tasks in the job, so that these requirements can be matched with agent capabilities.
%     Defines the artifacts that the build will produce.
%     Can only use artifacts produced in a previous stage.
%     Specifies any labels with which the build result or build artifacts will be tagged.
% 
% Task	
% 
%     Is a small discrete unit of work, such as source code checkout, executing a Maven goal, running a script, or parsing test results.
%     Is run sequentially within a job on a Bamboo working directory.
    
    
    
\subsection{کار مستند فنی}

ایجاد مستند فنی را می‌توان در قالب یک کار در این کارگزار در نظر گرفت. در این کار
ابتدا باید کد منبع از مخزن‌های کنترل نسخه دریافت شده و بر اساس آن مستند فنی
ایجاد شود. در نهایت مستند فنی ایجاد شده به صورت پرونده‌های فشرده ایجاد شود و به
عنوان مصنوعات در سیستم معرفی شود.


ایجاد پروژه

ایجاد نقشه

ایجاد وهله

ایجاد کار 



\subsubsection{کد منبع}

کد منبع از منبع کنترل نسخه دریافت می‌شود

\subsubsection{مستند فنی}

مستند فنی بر اساس آن ایجاد  می‌شود

\subsubsection{بسته بندی}

نتایج به دست آمده بسته بندی می شود